\section{讨论}

\subsection{Stage-3 surrogate 的筛查模型定位}

合成数据集由第~4.1 节中的降阶 \(q_1\) 拟合生成：原始 CDF 轨迹先经过清洗与时间对齐，再拟合至 \(q_1\) 形式，并在设定时间窗口内提供平滑外推目标。

Stage~1 从完整合成数据集中逐工况提取稳定的代表性均值轨迹，并验证了 $A$ 缩放假设。
但由于每个工况只保留一条代表性轨迹，单独使用 Stage~1 会忽略 plume-to-plume 与 shot-to-shot 变异，因此不适合作为最终筛查模型。

Stage~2 将监督扩展至完整合成数据集，并通过两个输出头同时学习外推窗口内的条件均值和条件离散度。
这也解释了它的 MAE 略高于 Stage~1：Stage~2 面对的是一个更困难、也更接近最终筛查需求的监督问题。

Stage~3 结合了前两个阶段的优势。
在原始观测较密集的早期时间段，它直接从测得的 CDF 数据学习；
在观测随时间推进而变得稀疏的晚期时间段，它按照第~2 章建立的右删失依据，通过 Gaussian KL 与 Stage~2 教师给出的条件均值和方差对齐。
由于教师和学生都输出单峰 Gaussian 分布，二者之间的 KL 散度具有闭式解，可直接由条件均值与方差计算，无需 Monte Carlo 采样，也不会额外引入采样噪声。
这种混合监督解释了 Stage~3 的早期保真度和更保守的晚期行为。

\subsection{已验证发现与适用边界}

在本文声明的范围内，有四项结论得到结果支持。
第一，光学工作流能够在归档尺度上从多孔 Mie 视频中提取可重复的穿深目标。
第二，降阶拟合同时提供轨迹平滑和面向删失观测的目标构造。
第三，分阶段 surrogate 从代表性均值预测推进到不确定度感知的完整数据集预测，并最终达到基于原始 CDF 精化的筛查性能。
第四，最终不确定度区间足够保守，可支持谨慎的早期工程排序，而不仅仅是点预测。

本文对结论适用范围作出明确限定。
所提取的喷雾边界被解释为特定光学协议下由 Mie 成像定义的液相包络观测量，而非完整的缸内液相真值 \cite{ecn_liquid_penetration,ecn_liquid_penetration_accessory}。
模型主要面向训练工况域内的插值预测，而非无约束外推；
它也不旨在替代 CFD、发动机试验或经验证的喷雾--壁面撞击物理模型，而是作为一层可审计的早期工程筛查工具，在启动更高保真验证前收窄设计空间。

目标定义本身也带有限制。
降阶拟合会收敛到以四分之一次方分支为主导的延拓，神经 surrogate 又通过正则化被约束为单调递增。
这些形状先验有助于获得符合直觉的轨迹增长，并修补观测窗口尾端的短时外推，但它们并不模拟 Mie 散射的成像机理。
本文进一步假设，在室温、非反应 Mie 实验中，晚期图像信号减弱主要来自液滴破碎、空间扩散、局部液滴数密度降低、散射强度跌入检测阈值，以及亮度倍数掩膜带来的检测副作用，而不是燃油在毫秒级窗口内快速蒸发消失。
因此，即使 plume 前沿在到达观测窗口边缘之前逐渐不可见，本文仍将其解释为液相几何包络可见性下降，而不是液滴运动已经物理终止。
这一解释与室温柴油蒸发相对缓慢的物理直觉一致，但本文没有同步液滴粒径、质量通量或蒸发率测量来直接验证它。
因此，得到的穿深应被理解为 Mie 定义的几何包络延续，而不是强液相喷雾尖端、局部液体质量、液滴尺寸、相态或壁膜形成潜力的直接证据 \cite{ecn_liquid_penetration,ecn_liquid_penetration_accessory}。

这一外推假设仍缺少直接对照验证。
原计划的验证路径是使用工业合作方提供的同类型喷嘴、相近工况和约两倍视场的宽视场单喷油束 Mie 视频，完整观测穿深轨迹直至平台区；
随后人为截断视场，仅用截断前时间序列拟合 \(q_1\) 模型，并将外推曲线与截断后仍可见的完整测量轨迹逐时刻比较。
该实验能直接检验本文 KD 链条的核心假设：参数形式 \(\sigma\!\left((t-t_0)/s\right)\cdot k_{1/4}\,t^{1/4}\) 在视场边界之外仍是可信延拓，而不是主要由形状先验支配的偏差。
由于合作方未能提供此类单喷油束数据，该验证未能在本文范围内完成；因此，使用 \(q_1\) 假设外推的可靠时间范围仍是一个开放经验问题。

上述边界对下游壁面筛查很重要。
喷雾与壁面的相互作用并不总是有害：如果大部分液相在到达壁面前已经蒸发，剩余的壁面相互作用可能危害更小，某些工况下甚至有助于卷吸和混合。
真正需要优先关注的是直接液滴撞壁和持续壁膜形成 \cite{MorieiraMotaPanao2010,liu2024marine_impingement,peraza2022swi,ahmad2025fueldilution}。
因此，本文模型更应被理解为对 Mie 定义喷雾包络可能接近壁面的保守筛查 surrogate，而不是对有害液态撞壁的相态分辨预测器。
液相包络几何与温度、蒸发、局部液滴质量通量及壁膜形成之间的对应关系并未在本文中建模或验证；
较高的 wall-interaction 分数只能表示几何接近风险较高，不能直接等同于有害液态撞壁或壁面润湿量。

下游发动机几何耦合同样是简化模型。
当前原型只用滑块--曲柄运动学平移活塞壁面 mask，没有将时变腔压、压力波动力学或大尺度气流运动反向纳入穿深预测。
这一简化继承自 OSCC 测量本身，因为每次喷射事件都发生在固定背压的刚性腔体中；
因此，surrogate 学到的是准静态舱内条件下的穿深，而不是完全耦合的缸内瞬态。
考虑到柴油喷雾穿深对环境气体密度敏感，这种解耦应被理解为务实建模选择，而不是压力和气流效应可忽略的证明 \cite{NaberSiebers1996}。
在本文目标场景中，即靠近上止点、时间尺度为毫秒级的短预喷窗口，这一近似仍可作为一阶工程筛查接受：当前玩具活塞设置下，\SI{5}{ms} 窗口内冠面位移小于 \SI{1}{mm}，而当前需求是几何排序，不是曲轴转角分辨的缸内真值预测。
当转速升高、喷射时刻远离上止点，或缸内瞬态压力与气流速度显著改变喷雾所见环境状态时，当前 wall-interaction 分数只能被读作解耦筛查 proxy。

另一个相关限制是喷雾锥角。
虽然图像处理阶段提取了 cone-angle proxy，但本文未在不同 campaign 间系统分析它，也没有将其提升为最终 surrogate 的学习目标。
因此，下游 wall-interaction 模型中的半锥角 \(\alpha\) 是固定输入，而不是经验证的时变观测量。
既有喷雾研究表明，扩张角会随瞬态发展、环境气体密度、喷射条件和喷嘴几何而改变 \cite{NaberSiebers1996,JungManinSkeenPickett2015}。
当前几何筛查只学习轴向穿深方向上的不确定度；由 \(\alpha\) 决定的横向包络仍是被假定的，而不是被预测的。

最后，结果还给出以下具体边界：
\begin{itemize}
	\item surrogate 仅在实验和合成数据共同支持的工况域内训练。
	      在真正分布外的压力、喷嘴几何、环境温度或喷射策略下，其预测不应被赋予同等可信度。

	\item 留一喷嘴族评估表明，分阶段 MLP 在未见喷嘴族上的 RMSE 仍低于两个物理基线，说明其均值预测具备一定跨族稳健性。
	      但最困难留出族出现低覆盖，且 Nozzle~3 的图像证据显示误差与先导液团--主射流反超、孔间非对称和阶跃式穿深有关。
	      这一现象与 split/multiple-injection 文献中报道的 wake-driven following-spray acceleration 和 catch-up/overtake motion 相一致：前一液团可通过尾流、气相轴向动量和局部混合场预扰动改变后一液团的穿透行为。
	      然而，本文没有同步气相速度场或液团身份追踪，因此只能将 Nozzle~3 解释为与该机制一致的图像证据，而不能单独证明其完整流体力学来源。
	      当前模型的主要跨族风险不是平均穿深整体失效，而是对未见喷嘴族中特殊早期瞬态的不确定度覆盖不足。

	\item 外部清洗评估报告了四次 anchor-off 重复训练，seed-to-seed RMSE 标准差为 \SI{0.120}{mm}。
	      该均值比单次点估计更可靠，但仍不等同于完整统计认证，因为尚未报告 bootstrap 置信带。

	\item 原始 held-out RMSE 和 MAE 基于 21{,}470 条轨迹中的 519{,}570 个逐点样本计算，但同一轨迹内的时间点误差彼此相关。
	      因此，独立样本数显著小于逐点样本数，任何基于逐点数量的置信表述都应结合这种相关性理解。

	\item 分阶段 surrogate 已在同一外部清洗协议下与校准 Hiroyasu--Arai、延迟 Naber--Siebers、稀疏异方差 Gaussian-process 基线以及自身 Stage~3 消融比较。
	      这些比较覆盖了物理相关式、非参数概率回归和本文监督策略的主要差异。
	      但 random forest、gradient boosting 或 quantile boosting 等轻量 tabular 回归器尚未系统纳入；
	      因此，本文不声称神经网络结构相对于所有确定性机器学习替代模型均具有最优边际收益。
\end{itemize}

\subsection{概率性碰壁筛查原型的未来验证路径}

第~4.7 节所描述的碰壁筛查原型实现了文献综述中主张的概率性 proxy，但逐帧碰撞概率评分 \(P_{\mathrm{hit}}(t)\) 尚未与实物喷雾--壁面观测比对。
这一空白是有明确定义的后续工作，而非开放性限制，因为所需参考数据已存在于工业合作方的实验数据集中。

最直接的验证目标是代表性活塞几何下的高速光学实验，即在发动机代表性边界条件下对喷射过程逐帧成像。
这类实验可直接观测喷雾在活塞冠附近的扩展情况，并与 \(P_{\mathrm{hit}}(t)\) 对比：
观测到喷雾接触的首帧给出碰撞起始的经验边界，喷射窗口内的累积喷雾暴露量则可与积分评分 \(E = \int P_{\mathrm{hit}}\,\mathrm{d}t\) 比较。
这一比较将在训练数据集无法提供的动态几何条件下检验轴向穿深 surrogate。

另一条互补路径是 CFD 数值模拟。
高保真喷雾计算能够分辨液相分布、壁面液膜质量和局部液滴浓度，而 Mie 成像定义的 surrogate 只用旋转 Gaussian 包络近似这些物理量。
该 surrogate 的定位并非取代 CFD，而是作为低成本上游筛查层：碰撞概率超过阈值的工况进入完整计算，安全低于阈值的工况则被过滤。
验证这一分层工作流的筛查召回率、假阴性率与计算量节省，需要一组 CFD 标注参考工况；当前数据集尚不具备这些。

两条验证路径都将解决当前原型中的结构性开放问题：固定半锥角 \(\alpha\) 是否可接受，或是否需要使用第~3 章提取的锥角 proxy 构建时变横向展宽模型。
只有获得逐帧匹配的参考数据后，对 \(\alpha\) 的灵敏度才能被可靠评估。
